\section{Introduction}

Large-area point clouds may contain substantially more points than a~model can process in a~single pass, while their spatial density varies because of acquisition geometry, sensor distance, and occlusions. Full-scene inference therefore requires a~pipeline that divides the data into bounded inputs and reconstructs predictions for the complete cloud~\citep{randlanet}.

This work was developed as part of the BRIK project, which included semantic segmentation of large point clouds collected near railway tracks. These clouds had non-uniform point density and numerous gaps. These conditions motivated the focus on density variation.

\subsection{Research Problem}

Partitioning and sampling determine the spatial extent, point distribution, and coverage presented to the model~\citep{randlanet,sensaturban}. Differences in sensors and acquisition conditions alter point-sampling patterns, and density variation is one component of the resulting domain gap~\citep{densityshift,complete-label,lidomaug}. The research problem is to determine how the external pipeline affects segmentation accuracy, execution time, and memory consumption particularly under density shift.

\subsection{Study Objective}

The objective is to design and evaluate an external pipeline for RandLA-Net against a~reference workflow on the BRIK project dataset and the DALES benchmark. The comparison covers data preparation, inference, and full-cloud reconstruction; the RandLA-Net architecture is held fixed.
